\documentclass[../../../include/open-logic-section]{subfiles}

\begin{document}
\olfileid[pt]{sth}{card-arithmetic}{expotough}

\olsection[Algumas Simplificações]{Algumas Simplificações com Exponenciação Cardinal}

Embora definir $\canonord$ fosse um pouco trabalhoso, o resultado é um
teorema útil sobre adição e multiplicação cardinais,
\olref[simp]{cardplustimesmax}. A exponenciação transfinita, porém,
não se simplifica tão diretamente. Para explicar por que, começamos
por um resultado que prolonga um padrão familiar do caso finitário
(embora a sua prova esteja a um nível elevado de abstração):

\begin{prop}\ollabel{simplecardexpo}
$\cardexpo{\cardfont{a}}{\cardfont{b} \cardplus \cardfont{c}} =
\cardexpo{\cardfont{a}}{\cardfont{b}} \cardtimes
\cardexpo{\cardfont{a}}{\cardfont{c}}$ e
$\cardexpo{(\cardexpo{\cardfont{a}}{\cardfont{b}})}{\cardfont{c}} =
\cardexpo{\cardfont{a}}{\cardfont{b} \cardtimes \cardfont{c}}$, para
quaisquer cardinais $\cardfont{a}, \cardfont{b}, \cardfont{c}$.
\end{prop}

\begin{proof}
Para a primeira afirmação, considere-se uma função $f \colon
(\cardfont{b}\disjointsum\cardfont{c}) \to \cardfont{a}$. Ora, ``parta-se'' isto, definindo $f_\cardfont{b}(\beta) = f(\beta, 0)$ para cada
$\beta \in \cardfont{b}$, e $f_\cardfont{c}(\gamma) = f(\gamma, 1)$
para cada $\gamma \in \cardfont{c}$. A aplicação $f \mapsto
(f_{\cardfont{b}} \times f_\cardfont{c})$ é !!a{bijection}
$\funfromto{\cardfont{b} \disjointsum \cardfont{c}}{\cardfont{a}} \to
(\funfromto{\cardfont{b}}{\cardfont{a}} \times
\funfromto{\cardfont{c}}{\cardfont{a}})$. 

Para a segunda afirmação, considere-se uma função $f \colon \cardfont{c} \to
(\funfromto{\cardfont{b}}{\cardfont{a}})$; assim, para cada $\gamma \in
\cardfont{c}$ temos alguma função $f(\gamma) \colon \cardfont{b} \to
\cardfont{a}$. Ora, defina-se $f^*(\beta, \gamma) = (f(\gamma))(\beta)$
para cada $\tuple{\beta, \gamma} \in \cardfont{b} \times \cardfont{c}$.
A aplicação $f \mapsto f^*$ é !!a{bijection}
$\funfromto{\cardfont{c}}{(\funfromto{\cardfont{b}}{\cardfont{a}})}
\to \funfromto{\cardfont{b} \cardtimes \cardfont{c}}{\cardfont{a}}$. 
\end{proof}

O que gostaríamos \emph{mesmo} era de uma maneira fácil de calcular
$\cardexpo{\cardfont{a}}{\cardfont{b}}$ quando tratamos de
cardinais infinitos. Eis um passo útil nessa direção:

\begin{prop}\ollabel{cardexpo2reduct}
Se $2 \leq \cardfont{a} \leq \cardfont{b}$ e $\cardfont{b}$ é
infinito, então $\cardexpo{\cardfont{a}}{\cardfont{b}} =
\cardexpo{2}{\cardfont{b}}$
\end{prop}

\begin{proof}
\begin{align*}
	\cardexpo{2}{\cardfont{b}} &\leq 
	\cardexpo{\cardfont{a}}{\cardfont{b}}\text{, pois $2 \leq \cardfont{a}$}\\
	&\leq \cardexpo{(2^\cardfont{a})}{\cardfont{b}}
	\text{, por \olref[opps]{lem:SizePowerset2Exp}}\\
	&= \cardexpo{2}{\cardfont{a} \cardtimes \cardfont{b}}
	\text{, por \olref{simplecardexpo}} \\
	&= \cardexpo{2}{\cardfont{b}}
	\text{, por \olref[simp]{cardplustimesmax}}
\end{align*}
\end{proof}

Não devemos esperar poder simplificar mais além disto,
pois $\cardfont{b} < \cardexpo{2}{\cardfont{b}}$ por
\olref[card-arithmetic][opps]{lem:SizePowerset2Exp}.
Contudo, isto não nos diz o que afirmar sobre
$\cardexpo{\cardfont{a}}{\cardfont{b}}$ quando $\cardfont{b} <
\cardfont{a}$. Naturalmente, se $\cardfont{b}$ é \emph{finito}, sabemos
o que fazer.

\begin{prop}
Se $\cardfont{a}$ é infinito e $n \in \omega$, então
$\cardexpo{\cardfont{a}}{n} = \cardfont{a}$
\end{prop}

\begin{proof}
$\cardexpo{\cardfont{a}}{n} = \cardfont{a} \cardtimes \cardfont{a}
\cardtimes \ldots \cardtimes \cardfont{a} = \cardfont{a}$, por \olref[simp]{cardplustimesmax}.
\end{proof}
\noindent 
Adicionalmente, nalguns outros casos, podemos controlar o tamanho de
$\cardexpo{\cardfont{a}}{\cardfont{b}}$:

\begin{prop}
Se $2 \leq \cardfont{b} < \cardfont{a} \leq
\cardexpo{2}{\cardfont{b}}$ e $\cardfont{b}$ é infinito, então
$\cardexpo{\cardfont{a}}{\cardfont{b}} = \cardexpo{2}{\cardfont{b}}$
\end{prop}

\begin{proof}
$\cardexpo{2}{\cardfont{b}}\leq \cardexpo{\cardfont{a}}{\cardfont{b}}
\leq \cardexpo{(\cardexpo{2}{\cardfont{b}})}{\cardfont{b}} =
\cardexpo{2}{\cardfont{b}\cardtimes\cardfont{b}} =
\cardexpo{2}{\cardfont{b}}$, raciocinando como na \olref{cardexpo2reduct}.
\end{proof}
\noindent 
Para além deste ponto, as coisas tornam-se bastante mais subtis.

\end{document}
